%% notes-tex/019-simb-multimodal/sections/5-amino-acid.tex
%% [[experiments.023-metabolome-betaxanthin-joint.scripts.betaxanthin_amino_acid_predictivity]]
%% SOURCE: coupling numbers from
%% experiments/023-metabolome-betaxanthin-joint/scripts/betaxanthin_amino_acid_predictivity.py
%% -> results/betaxanthin_amino_acid_predictivity.json. Scores from
%% experiments/019-simb-multimodal/results/round_leaderboards.csv. Convergence behavior from
%% experiments/022-mulleder-metabolome/results/training_length_analysis.json.

\section{Amino-acid pools, and whether they help production\secstatus{todo}}
\label{sec:amino-acid}

\subsection{The target\secstatus{todo}}

Mulleder 2016 (doi:10.1016/j.cell.2016.09.007) measured the intracellular concentration of
19 amino acids in 4,678 single-deletion strains. There is one replicate per strain and no
released standard error, so the dataset admits \textbf{no within-dataset reproducibility
ceiling}. Every correlation involving it is attenuated by an unmeasured amount, and that
qualification travels with every number below.

The best validation \file{pearson_per_feature} is $\mathbf{0.209}$, run
\texttt{24vlgope} in \file{torchcell_022_mulleder19_v4}, peaking at epoch 402 of 999.
The median across all 97 runs in the strand is $0.0570$, so the leaderboard is a long tail
under one good configuration rather than a plateau. Earlier rounds read $\le 0.025$, then
$0.171$, so the trajectory is upward across rounds.

\notesec{That median is corrected, and the correction matters for every median in this
document} Revision 2 reported $0.0498$ ``across the 31 runs''. Those 31 were not a sample of
the strand: they were the subset whose validation curves had been fetched, and the fetch
rule selects the top runs per project by final score and by epoch count
(\cref{sec:coverage}). A maximum-finder is the right instrument for a strand maximum and the
wrong one for a median, since the runs it skips are systematically the weak ones. Reading
all 97 moves the median up to $0.0570$. Maxima are unaffected by construction.

\subsection{These runs turn within their budget, and expression does not\secstatus{todo}}

Twenty-four runs in the preceding round were scored for their convergence shape. The
typical run peaks at epoch 48 to 134, continues for another 41 epochs with validation slope
turning negative while training slope stays positive, and is classified converged then
overfitting. None of them hit the epoch cap.
\src{experiments/022-mulleder-metabolome/results/training_length_analysis.json}

\notesec{Said carefully, because ``converged'' has already misled us once on this project}
What is observed is that the validation metric reaches a maximum inside the budget and then
declines, which is what makes these runs scorable. That is weaker than convergence, and the
distinction is not pedantic: the expression runs were also called converged, on curves that
dipped early and flattened, and only at several thousand epochs did it become clear that
Pearson was still climbing and squared error would come back down to a late minimum. The
same could in principle be true here beyond epoch 999, and no amino-acid run has been
extended to find out.

What can be said is the comparison, which does not depend on the word. Amino-acid runs reach
a maximum well inside a 999-epoch budget; expression runs have not reached one at 10,000.
The two use the same architecture and the same encoder, so whatever makes expression climb
for thousands of epochs is a property of the expression target or its loss rather than of the
model.

\subsection{The precursor hypothesis is wrong, and the profile hypothesis is right\secstatus{todo}}

Betalains derive from L-tyrosine (tyrosine to L-DOPA to betalamic acid), and a betaxanthin
is betalamic acid condensed with an amino acid. The mechanistic prediction is that a
deletion's free amino-acid pool carries information about how much betaxanthin that same
deletion makes. Both screens are keyed by systematic ORF and share 4,432 deletions, so the
prediction is directly testable with no model in the loop.

\notesec{What the prediction actually is, since it is easy to read as a model result and
is not one} No neural network appears anywhere in this subsection. The procedure is:

\begin{enumerate}
\item Join the two screens on systematic ORF. One row per deletion: 19 measured amino-acid
  concentrations from Mulleder, and one betaxanthin value from Cachera.
\item Take $\log$ of each concentration (they are skewed positive quantities) and
  $z$-score each of the 19 columns.
\item Fit \textbf{ridge regression} from those 19 numbers to the betaxanthin value, with
  the penalty chosen by inner cross-validation.
\item Score \textbf{out of fold}: five-fold cross-validation, so every prediction is made
  by a model that never saw that deletion, then correlate the held-out predictions against
  the measurements. Repeat over three fold shuffles and report the mean.
\end{enumerate}

So $r = 0.298$ means: \emph{knowing a deletion's 19 amino-acid concentrations lets you
predict its betaxanthin at $r = 0.298$ on deletions you did not fit on}. The comparison row
``tyrosine alone'' is the same procedure with one column instead of nineteen. It is a
statement about the two measurements, and it puts a floor under what any model sharing a
representation between these phenotypes could hope to exploit.

\begin{table}[htbp]
\centering\small
%% SOURCE: results/betaxanthin_amino_acid_predictivity.json (cross_validated_uncertainty_key)
\begin{tabular}{lrrrr}
\toprule
predictor of betaxanthin & $r$ & SE at this $n$ & SD across shuffles & $n$ \\
\midrule
tyrosine alone (the precursor)         & 0.0641 & 0.0150 & 0.0039 & 4{,}432 \\
\textbf{19 amino acids jointly}        & \textbf{0.2980} & 0.0137 & 0.0029 & 4{,}432 \\
\midrule
\multicolumn{5}{l}{\emph{on the 3{,}708 deletions with a Costanzo fitness record}} \\
single-mutant fitness alone            & 0.1508 & 0.0161 & 0.0028 & 3{,}708 \\
19 amino acids                         & 0.1763 & 0.0159 & 0.0025 & 3{,}708 \\
19 amino acids and fitness             & 0.2001 & 0.0158 & 0.0010 & 3{,}708 \\
19 amino acids, fitness regressed out of both & 0.1326 & 0.0161 & 0.0026 & 3{,}708 \\
\midrule
\multicolumn{5}{l}{\emph{on the 724 with no Costanzo record}} \\
19 amino acids                         & \textbf{0.4547} & 0.0295 & 0.0219 & 724 \\
\bottomrule
\end{tabular}
\caption[]{Cross-validated ridge, five folds, three shuffles; amino acids enter as $\log$
concentration, $z$-scored. \textbf{Two different spreads are reported because they answer
different questions and quoting only one is misleading.} The \emph{SE at this $n$} is the
Fisher standard error, $1/\sqrt{n-3}$ mapped back at the observed $r$, and it is the one to
use for ``how precisely is this correlation measured'': it is what separates the $0.298$
from the $0.064$, at $17$ standard errors. The \emph{SD across shuffles} is fold-assignment
noise on the same deletions, which is why it is an order of magnitude smaller; it says the
procedure is stable, not that the estimate is precise. A third spread, the SD across the
five folds, runs $0.025$ to $0.064$ and is an upper reference rather than a standard error,
since each fold estimates on $n/5$ deletions.}
\label{tab:aa-predictivity}
\end{table}

\notesec{Tyrosine is not the axis} Its marginal correlation with betaxanthin is
$r = -0.076$ ($p = 3.9\times10^{-7}$), which is negative, and it ranks sixth of nineteen by
absolute correlation. The strongest is methionine at $+0.140$, followed by aspartate
$+0.127$ and arginine $-0.118$. No amino acid exceeds $|r| = 0.15$. Correcting for the
betaxanthin side's known reliability of $0.836$ raises these by a factor of $1.094$, which
does not change the reading; the Mulleder side's correction is unmeasurable and would raise
them further by an unknown amount.

\notesec{The profile is the axis} The 19-dimensional pool predicts betaxanthin at
out-of-fold $r = 0.298$, roughly $4.7\times$ the precursor alone.

Revision 2 set that beside the CGT's own genotype-to-betaxanthin score and said a linear map
``recovers most of what the model recovers from genotype''. \textbf{That comparison was not
valid} and it is withdrawn: the two numbers use different splits and different scoring
rules, and one takes measured metabolites as inputs while the other takes only the genotype.
The question the comparison was reaching for is answered properly below, by a measurement
rather than by placing two numbers next to each other.

\subsection{Is the coupling redundant with what betaxanthin supervision already teaches\secstatus{todo}}
\label{sec:aa-redundancy}

This is the objection that decides whether the coupling is worth building on, and revision 2
did not address it. Both the amino-acid pool and betaxanthin are readouts of the same
deletion. So an $r$ of $0.298$ is equally consistent with two very different worlds:

\begin{enumerate}
\item[(i)] the pool carries betaxanthin-relevant information that a genotype-to-betaxanthin
  model does \emph{not} already extract, in which case a shared encoder has something to
  gain; or
\item[(ii)] both are shadows of one genotype axis the Cachera labels already teach, in which
  case no amount of architecture can beat betaxanthin-only training, and the $0.298$ is
  real but useless.
\end{enumerate}

\notesec{The measurement that separates them} Take a trained betaxanthin-only model, read
its \emph{held-out} predictions, and regress those out of both sides before asking what the
amino-acid pool still predicts. If world (ii) holds, nothing survives. Run over the ten
grid cells that dumped test predictions, on the 816 held-out deletions carrying both
measurements:

\begin{table}[htbp]
\centering\small
%% SOURCE: experiments/023-metabolome-betaxanthin-joint/scripts/bx_aa_incremental_over_model.py
%% -> results/bx_aa_incremental_over_model.json
\begin{tabular}{lrrrr}
\toprule
control model (grid cell) & $r$ of the model & $r$ of 19 AA & 19 AA given the model & corr of the two \\
\midrule
\file{s08_L6_maskon_lr0.0001} & $0.3564$ & $0.2720$ & $\mathbf{0.2179}$ & $0.2185$ \\
\file{s09_L6_maskon_lr0.0001} & $0.3305$ & $0.2720$ & $0.2262$ & $0.2050$ \\
\file{s01_L2_maskon_lr0.0001} & $0.3074$ & $0.2720$ & $0.2139$ & $0.2263$ \\
\file{s05_L2_maskoff_lr0.0001} & $0.2913$ & $0.2720$ & $0.2259$ & $0.1768$ \\
\file{s04_L2_maskoff_lr0.0001} & $0.2761$ & $0.2720$ & $0.2380$ & $0.1660$ \\
\bottomrule
\end{tabular}
\caption[]{Against the strongest available control, the amino-acid pool keeps $0.2179$ of
its $0.2720$, which is $80\,\%$, after the betaxanthin model's own held-out prediction is
removed from both sides. At $n = 816$ that is $6.5$ Fisher standard errors from zero. The
last column is the direct check on the same point: the model's predictions and the
pool's predictions correlate at only $0.22$, where near-total redundancy would need
something above $0.75$. Five further collapsed cells with $r_{\text{model}} \le 0.10$ are
omitted, since residualizing on a failed model removes nothing by construction.}
\label{tab:aa-incremental}
\end{table}

\notesec{So world (i) holds} There is betaxanthin-relevant information in the free
amino-acid pool that the genotype-to-betaxanthin map does not capture, and the joint arm's
failure to use it is therefore a statement about the auxiliary head rather than about the
signal.

\notesec{Three things this does not establish, kept separate because they are separate}
It controls for a \emph{linear} function of the model's prediction, so non-linear redundancy
is untouched. The strongest control available is $r = 0.3564$, below the strand's best
score, so redundancy against a better model than any that dumped predictions is not ruled
out. And ``the coupling is real and incremental'' is a different claim from ``an auxiliary
head is the mechanism that converts it into a gain''; the second is what the next subsection
finds against.

\notesec{Why fitness enters at all, since it is a third dataset and looks like a detour}
Both screens are deletion phenotypes and both respond to how well the strain grows. A slow
deletion can have a distorted amino-acid pool \emph{and} a distorted pigment readout with
no metabolic relationship between the two, which would produce exactly the correlation
above for an uninteresting reason. Costanzo single-mutant fitness (KanMX deletions, 30 C)
is the standard measurement of that shared nuisance, so it is regressed out of both sides
and the fit recomputed on the residuals. It is a control, not a third phenotype in the
story.

Single-mutant fitness alone predicts betaxanthin at $0.151$, so the shared growth component
is real and substantial. Regressing it out of both sides leaves the amino-acid profile at
$0.133$ against $0.176$ without the control on the same genes, so roughly three quarters of
its predictive power survives. Combining the two reaches $0.200$, above either alone.

\notesec{Correction: most of that apparent shrinkage was the gene set, not the control}
Requiring a Costanzo record shrinks the population from 4,432 to 3,708, and the 724 it drops
are not a random quarter. Rather than accept that, every single-deletion fitness source in
the repository was pooled. Kuzmin does \emph{not} help, and the reason is structural: Kuzmin
2018 and 2020 screened the same SGA deletion array Costanzo did, agreeing with it at
$r = 0.947$ to $0.995$ on the overlap, so between them they recover only 56 of the 724.
Baryshnikova 2010 recovers 375, O'Duibhir 2014 a further 167, and the Costanzo NatMX query
arm, which the incumbent control excluded, another 317. The union controls \textbf{4,214 of
the 4,432}, leaving 218.

\begin{table}[htbp]
\centering\small
%% SOURCE: experiments/023-metabolome-betaxanthin-joint/scripts/
%% smf_coverage_fitness_control_sources.py -> results/
\begin{tabular}{lrrrr}
\toprule
gene set & $n$ & 19 AA, no control & fitness alone & 19 AA, fitness controlled \\
\midrule
Costanzo KanMX only (incumbent) & 3{,}708 & $0.1763 \pm 0.0025$ & $0.1508$ & $0.1326 \pm 0.0026$ \\
\textbf{union of all SMF sources} & \textbf{4{,}214} & $\mathbf{0.2884 \pm 0.0014}$ & $0.1332$ & $\mathbf{0.2699 \pm 0.0017}$ \\
\midrule
all shared, no control possible & 4{,}432 & $0.2980 \pm 0.0029$ & -- & -- \\
still uncontrolled after the union & 218 & $0.3150 \pm 0.0312$ & -- & -- \\
\bottomrule
\end{tabular}
\caption[]{The fitness control costs far less than it appeared to. On the incumbent gene set
the control looks like it removes $0.044$ of predictive power ($0.176 \to 0.133$); on the
enlarged set it removes $0.019$ ($0.288 \to 0.270$). Most of the apparent shrinkage was the
population that requiring a Costanzo record selects, not the control itself. Donor screens
are mapped onto the Costanzo scale by least squares on their overlap and standardized within
donor; transfer correlations run $0.951$ (NatMX) and $0.891$ (Baryshnikova) down to $0.511$
(Kuzmin 2020), so the weakest donors carry the most transfer error. The incumbent row
reproduces the committed result to 13 decimal places, which is what makes the two rows
comparable.}
\label{tab:aa-fitness-union}
\end{table}

So the conclusion strengthens rather than weakens: \textbf{$94\,\%$ of the amino-acid
profile's predictive power survives controlling for single-mutant fitness} on the population
where the control can be applied at all.

\notesec{And the deletions no fitness screen reaches carry the extremes} The 724 deletions
with no Costanzo KanMX record have a betaxanthin standard deviation of $0.914$, against
$0.464$ among the 3,708 that do have one, so they are roughly twice as spread out in
production and carry a disproportionate share of the extreme producers. On that same 724 the
amino-acid profile predicts at $0.455$. Even after the union above, 218 remain uncontrolled,
and the profile predicts on those at $0.315$.

\emph{Why they are missing is itself informative, and unverified here.} A gene absent from a
KanMX deletion-collection fitness screen is typically one whose deletion grows too poorly to
score, or that the collection never carried. Either way the missing set is enriched for
strongly-affected strains, which is the plausible reason it carries the extremes. That
mechanism has not been tested.

\tcfigfit{figures/betaxanthin_amino_acid_predictivity.pdf}{(a) Marginal correlation of each
amino acid with betaxanthin over the 4,432 shared deletions; tyrosine highlighted. These
reproduce \file{experiments/019-simb-multimodal/results/pigment_noise_ceiling.json}
exactly, which the generating script asserts rather than assumes. (b) Cross-validated ridge
fits from \cref{tab:aa-predictivity}. (c) Median tyrosine by betaxanthin decile. The
relationship is non-monotone: tyrosine is highest in the lowest-production decile, falls
through the middle, and rises again at the top, which is why a linear correlation reads
near zero. The axis spans $2.65$ to $2.89$ mM, so the swing is about $9\,\%$ of the
median.}{fig:aa-coupling}

\notesec{A mechanism, labeled as unverified} The Cachera background carries
\gene{ARO4}$^{\mathrm{K229L}}$ and \gene{ARO7}$^{\mathrm{G141S}}$, feedback-resistant
alleles that release the shikimate pathway from the control that shapes the pools Mulleder
measured in the plain deletion collection. A plausible reading is that free tyrosine in an
unengineered strain is not the quantity that limits flux in an engineered one. This has not
been tested and no experiment here can test it.

\subsection{Does the metabolome head actually help the betaxanthin head\secstatus{todo}}

%% PLACEHOLDER-023: paired arm analysis, filled from round_leaderboards.csv
The claim under test is narrow and was stated before the round: in Delta grid cell
\file{s04_L6_maskon_lr0.0001}, the joint arm read $+0.0203$ over the control on
\file{val/betaxanthin/pearson_per_feature} at a matched budget of 333 epochs. That was
one paired difference at one seed with no estimate of the noise floor on it. The same grid
read $-0.049$ and $-0.060$ in the two $L = 2$ cells, both run to 1,000 epochs, so the sign
flips with depth.

%% GENERATED by experiments/019-simb-multimodal/scripts/build_retrospective_tables.py
%% SOURCE: results/round_leaderboards.csv, strand betaxanthin_amino_acid_joint.
\begin{table}[htbp]
\centering\small
\begin{tabular}{lrrrrl}
\toprule
grid cell & control & joint & $\Delta$ & aux $r$ & epochs (ctrl/joint) \\
\midrule
\texttt{s00\_L2\_maskon\_lr0.0001} & nan & nan & +nan & -- & nan/nan \\
\texttt{s02\_L2\_maskoff\_lr0.0001} & nan & nan & +nan & -- & nan/nan \\
\texttt{s04\_L6\_maskon\_lr0.0001} & nan & nan & +nan & -- & nan/nan \\
\texttt{s05\_L6\_maskon\_lr0.001} & nan & nan & +nan & -- & nan/nan \\
\texttt{s06\_L6\_maskoff\_lr0.0001} & 0.4261 & nan & +nan & -- & 998/nan \\
\texttt{s07\_L6\_maskoff\_lr0.001} & nan & nan & +nan & -- & nan/nan \\
\texttt{s08\_L4\_maskon\_lr0.0001}$^{\dagger}$ & 0.4149 & 0.4193 & \textbf{+0.0043} & 0.114 & 998/704 \\
\texttt{s09\_L4\_maskon\_lr0.001} & nan & nan & +nan & -- & nan/nan \\
\texttt{s10\_L4\_maskoff\_lr0.0001} & nan & nan & +nan & -- & nan/nan \\
\texttt{s11\_L4\_maskoff\_lr0.001} & nan & nan & +nan & -- & nan/nan \\
\texttt{s01\_L2\_maskon\_lr0.001}$^{\ast}$ & 0.0287 & 0.0413 & \textbf{+0.0125} & 0.026 & 998/998 \\
\texttt{s03\_L2\_maskoff\_lr0.001}$^{\ast}$ & 0.0200 & nan & +nan & -- & 998/nan \\
\midrule
mean, cells where the control learned (9 cells) & & & +nan $\pm$ nan & & 0 of 9 positive \\
mean, all matched-exposure cells (11 cells) & & & +0.0125 $\pm$ nan & & 1 of 11 positive \\
\bottomrule
\end{tabular}
\caption[]{Betaxanthin score with and without a 19-amino-acid auxiliary head, paired within grid cell, on \texttt{val/betaxanthin/pearson\_per\_feature}. $\Delta$ is joint minus control. \emph{aux} is the auxiliary head's own score, carried because an auxiliary head at $r \approx 0$ that still moved the primary metric would mean the gain came from regularization rather than from shared metabolic signal, and those are different findings. $^{\dagger}$ marks a cell whose two arms ran epoch counts differing by more than 50; since the score is a running maximum, such a difference is part effect and part unequal exposure, and those rows are excluded from both means. $^{\ast}$ marks a cell whose control arm scored below 0.10, meaning neither arm learned the task; every one of them is an $\mathrm{lr} = 10^{-3}$ cell, and a difference between two failed runs measures nothing about the auxiliary head.}
\label{tab:bx-aa-paired}
\end{table}


Scored uniformly across the grid, the effect is \textbf{negative and small}, and the
honest summary of it is three sentences rather than three readings.

The mean over the cells whose control arm learned the task at all is negative, the spread
across cells is several times the effect, and every cell carries $n = 1$ per arm except
\texttt{s08}, so there is no within-cell noise floor to measure the difference against.
Depth and the graph mask each move the result by more than the auxiliary head does, which
means any single-cell claim about the head is smaller than the variation across the cells it
could have been made in. And the auxiliary head is itself weak wherever the primary head is
good, scoring well under the $0.209$ the amino-acid strand reaches when it is the only head,
so the joint arm is paying capacity for a head that is underperforming.

\notesec{Read this as one crude attempt, not as a verdict} Adding an unstructured second
head is the least informative version of the question. Nothing in the architecture says the
two phenotypes are connected through a pathway, so the encoder is being asked to discover a
stoichiometric relationship from 19 noisy concentrations and one pigment readout. Given
\cref{tab:aa-incremental}, which shows the information is genuinely there and genuinely
incremental, the finding to carry forward is that \textbf{the route by which a second head's
representation reaches the first head's readout is the thing that failed}, and that a
constraint-based coupling is the version worth trying (\cref{sec:directions}).

The designed replication,
\file{experiments/023-metabolome-betaxanthin-joint/conf/igb_bx_aa.yaml}, adds a third
setting: an auxiliary head restricted to \textbf{tyrosine only}, against a common control,
on a gene set pinned by \file{require_head_targets} so no arm can win by seeing more rows.
\textbf{It has not been run.} Given \cref{tab:aa-predictivity} its tyrosine arm is now
predicted to be the null one and its 19-amino-acid arm the live one, which is a sharper
prediction than the config was written against.

\begin{keybox}
\textbf{The coupling is real, incremental, and not being used.} The free amino-acid profile
of a deletion predicts that deletion's betaxanthin at $r = 0.298$ out of fold, against
$0.064$ for tyrosine, the chromophore precursor, so the axis is the profile and not the
precursor. It is not a growth artifact: pooling every single-deletion fitness screen in the
repository controls 4,214 of the 4,432 shared deletions, and $94\,\%$ of the profile's power
survives that control. It is not redundant with betaxanthin supervision either: $80\,\%$ of
it survives residualizing out a trained betaxanthin model's own held-out predictions, which
correlate with the profile's predictions at only $0.22$. And an unstructured auxiliary head
does not convert any of that into a gain, scoring negative across the grid. \textbf{The
bottleneck is the route from a second head's representation to the first head's readout, not
the availability of the signal.}
\end{keybox}

\notesec{Cannot say} How much of the $0.298$ is attenuated by Mulleder's unmeasurable noise,
so it is a floor rather than an estimate. Whether the redundancy result survives a
\emph{non-linear} control, or a control model stronger than the $0.356$ best among those
that dumped predictions. And whether any architecture converts the coupling into a gain,
since the only mechanism tried is the one least likely to work.
